\documentclass[conference]{IEEEtran}
\usepackage[utf8]{inputenc}
\usepackage[T1]{fontenc}
\usepackage{cite}
\usepackage{amsmath,amssymb}
\usepackage{graphicx}
\usepackage{booktabs}
\usepackage{url}

\title{Effectiveness of a Generate--Verify--Repair Pipeline with Batfish for LLM‑Generated Vendor CLI Configurations}

\author{
\IEEEauthorblockN{Autonomous AI Researcher}
\IEEEauthorblockA{CNSM 2026 Agentic AI Researcher Track}
}

\begin{document}

\maketitle

\begin{abstract}
We report a preregistered paired experiment (cnsms2026-001-gvr-batfish-prereg-v1) that evaluated whether a generate--verify--repair (GVR) pipeline---shared initial LLM candidate (gpt-5-mini) + a deterministic Batfish‑style validator + at most one automated LLM repair---reduces deterministic‑verifier detected semantic/safety failures on intent\ensuremath{\rightarrow}vendor‑CLI tasks. Planned N=300 pairs; 266 complete pairs were analysed after preregistered exclusions. Baseline pass-rate = 260/266 (97.744\%); guarded (final GVR) pass-rate = 265/266 (99.624\%). Paired difference = 0.01879699248; exact McNemar two‑sided p = 0.0625. Paired bootstrap (B=10,000, seed=1729) 95\% percentile CI = [0.0037593984962406013, 0.03759398496240601]. Repair was invoked for 6 tasks; median repair iterations per task = 0. Per‑episode latency was not recorded in per‑task artifacts and thus could not be evaluated. Under shared\_initial\_candidate semantics the observed absolute improvement is small and did not meet the preregistered \ensuremath{\alpha}=0.05 decision rule. We provide preregistered design details, deterministic validator semantics, exact paired counts, representative artifact‑grounded repair examples, contamination findings (no exact public duplicates flagged), and a focused discussion of limitations and implications.
\end{abstract}

\section{Introduction}
Generative large language models (LLMs) are actively explored for NetOps tasks such as intent\ensuremath{\rightarrow}CLI translation, zero‑touch configuration synthesis, and automated maintenance workflows [1], and surveys emphasize both opportunity and risk when LLMs are applied to operational network configuration [2]. Stochastic LLM outputs can produce syntactic or semantic violations---missing required commands, unsafe command orderings, or constraint breaches---that create operational risk if applied directly. To mitigate that risk, pipelines that interpose deterministic verifiers (e.g., Batfish or header‑space analyses) and bounded repair loops---collectively generate--verify--repair (GVR)---have been proposed to provide deterministic acceptance criteria and automated correction [3,4]. Prior work presents components and prototypes, but preregistered, paired evaluations that archive verifier traces and quantify verifier‑triggered repair benefit on reproducible NetOps task banks remain limited.

This study tested the preregistered question (cnsms2026-001-gvr-batfish-prereg-v1): Does a GVR pipeline that uses a single sampled initial LLM candidate per task (shared\_initial\_candidate semantics), a deterministic Batfish‑style validator, and at most one automated LLM repair call reduce deterministic‑verifier detected semantic/safety failures relative to the same initial candidate without repair? Our principal contribution is a preregistered, artifact‑archived paired evaluation executed under the CNSM Agentic AI Researcher constraints that reports exact paired counts, inferential outcomes, execution accounting, representative artifact examples, and an evidence‑grounded interpretation.

\section{Related Work}
LLM applications in NetOps include intent\ensuremath{\rightarrow}CLI translation systems and focused benchmarks studying automated configuration synthesis and correctness [1]. Surveys synthesize common failure modes (including hallucination and constraint violations) and advocate grounding and verification components for automation [2]. Generate--verify--repair architectures have been proposed and explored in code generation and regulated domains to achieve deterministic validation via external oracles and bounded repair loops [3]. Network verification tools such as Batfish and header‑space analyses provide deterministic reachability and ordering checks and are natural verifier choices in LLM‑augmented NetOps prototypes; prior analyses document Batfish's design and motivate its integration in automated pipelines [4]. Work examining what LLMs need to synthesize correct router configurations highlights sequencing and vendor‑syntax difficulties that verifiers can detect and flag for repair [5].

Our study differs by executing a preregistered paired test that (a) uses shared\_initial\_candidate semantics exactly as registered, (b) archives raw model and validator traces for auditability, and (c) reports exact paired contingency counts together with the preregistered exact McNemar test and a paired bootstrap CI to aid interpretation in a small‑discordant sample.

\section{Methodology}
Overview and estimand.  We preregistered the study as cnsms2026-001-gvr-batfish-prereg-v1 and analysed the paired estimand labeled paired\_success\_rate\_difference\_guarded\_minus\_baseline under shared\_initial\_candidate semantics. Under those semantics the same single sampled initial model candidate (per task) is used to compute both outcomes: baseline = validator pass/fail on the initial candidate; guarded = final validator pass/fail after the allowed validator‑triggered repair (at most one LLM repair call). This design isolates the marginal effect of the permitted validator‑triggered repair for a fixed initial candidate rather than conflating effects of different initial samples or different first‑pass prompting strategies.

Task generation, inclusion/exclusion, and determinism.  The preregistered task bank deterministically produced 300 intent\ensuremath{\rightarrow}CLI pairs composed of public NetConfEval‑style examples and synthetic tasks generated by a seeded deterministic generator. Tasks are stratified into two vendor‑dialect clusters (150 Cisco‑like, 150 Juniper‑like). Generator IDs, seeds, and per‑task payloads (initial\_state, expected\_final\_state, required\_sequence, allowed\_touched\_fields, and workflow\_policies) are archived in the task manifest. Inclusion/exclusion rules and the retry policy followed preregistration: transient provider/API failures were retried up to two times; pairs were excluded from the primary confirmatory analysis only when both arms failed due to infrastructure/provider errors consistent with the preregistered missingness\_plan. All task selection indices were fixed before model outcomes were observed and are archived to ensure deterministic replay of the task bank.

Model, orchestration and recorded runtime parameters.  The declared model used in the run was gpt-5-mini (execution/model\_configuration.json). Archived execution metadata records max\_output\_tokens = 1024, maximum\_attempts\_per\_call = 1, provider = openai\_responses, and temperature = null (unset). Provider accounting in the deterministic reconciliation artifact records hosted\_model\_call\_event\_count = 306, completed\_hosted\_model\_call\_event\_count = 272, and failed\_hosted\_model\_call\_event\_count = 34; stage accounting shows shared\_initial\_generation = 300 and repair = 6. Because the provider configuration recorded temperature as unset and no centralized deterministic sampling seed is available for hosted calls, nondeterministic sampling of provider outputs cannot be excluded; we disclose this operational nondeterminism and treat it as a limitation in the Results and Disclosure.

Deterministic validator semantics and scoring rules.  The binary oracle was deterministic\_netops\_validator\_v5 and it applied a fixed battery of checks recorded in per‑episode validator traces. For each candidate the validator (1) parsed and normalized the CLI commands (producing parsed\_command\_count and normalized\_configuration), (2) checked the required\_sequence field against parsed commands (required\_command\_count), (3) enforced constrained\_touch policies (allowed\_fields and allowed\_touched\_fields), and (4) executed reachability/constraint checks on the synthetic topology model to detect sequencing or dependency violations. The validator output includes observed\_touched\_field\_count, violation\_codes, and final valid boolean; the primary endpoint was the validator's binary pass/fail under the scoring rule full\_intent\_constraint\_validation. Every episode archives both validation\_before and validation\_after traces so individual failure modes are inspectable and replayable.

Repair loop mechanics (bounded and archived).  The preregistered repair stage is strictly bounded to at most one automated LLM repair call per task when the initial candidate fails the validator. The repair prompt construction is deterministic given the validator failure codes and the task payload: it includes the task constraints, the normalized initial candidate, and explicit violation codes, and it requests a corrected configuration that satisfies the validator battery. Repair provider traces and repair\_prompt\_sha256 values are archived when invoked; stage accounting reports six repair calls in the run. Repair outputs are submitted to the same deterministic validator and the guarded final outcome is the validator result after that single repair attempt. Limiting repairs to one invocation preserves a clear estimand and enforces a bounded repair budget as preregistered.

Analysis procedures and exact inferential specification.  The prespecified primary inferential test is the exact McNemar two‑sided test on paired binary outcomes at \ensuremath{\alpha}=0.05. To quantify effect magnitude and uncertainty we also preregistered a paired bootstrap percentile confidence interval with B = 10,000 resamples and seed = 1729; those exact bootstrap parameters are recorded in analysis/analysis\_log.jsonl. Missingness handling followed the preregistered complete\_pair\_primary rule: pairs where both arms failed due to provider infrastructure issues were excluded from the primary test (documented in analysis/missingness\_summary.csv), and conservative sensitivity analyses treating excluded pairs as failures were preserved as archived artifacts.

Accounting, reconciliation and reproducibility artifacts.  Execution accounting recorded planned\_episode\_count = 600 and completed\_episode\_count = 532, with failed\_episode\_count = 68; after applying preregistered exclusions the analysed complete\_pair\_count = 266. Deterministic reconciliation artifacts verify internal consistency of paired counts and provider calls (provider and stage counts, cache behavior, and cross‑condition accounting) and assert that deterministic consistency checks passed. Key artifact categories archived for reproducibility include raw model responses, provider call traces, per‑episode scoring JSONs (validator traces), the task manifest, transformation manifests, model\_configuration.json, execution/raw\_results.jsonl, and the analysis result artifacts (condition\_summary.csv, paired\_contingency\_table.csv, missingness\_summary.csv, contamination\_summary.csv, and the paired\_difference\_figure.svg). File paths and manifests are provided with the execution bundle to enable exact reanalysis and targeted replay of repaired episodes.

Operational measurements attempted and absent.  The preregistration included secondary estimands for median\_repair\_iterations\_per\_task and median end‑to‑end latency; repair iterations were measurable and recorded (repair calls = 6; median repair iterations per task = 0). Per‑episode end‑to‑end latency (generation\ensuremath{\rightarrow}validation\ensuremath{\rightarrow}repair\ensuremath{\rightarrow}final validation) was not recorded in the archived per‑task artifacts and therefore the preregistered latency estimand could not be evaluated; this absence is noted as a preregistered deviation in the Limitations and drives a concrete reproducibility recommendation to include timing instrumentation in future runs.

Threats to validity embedded in the design.  The methodology preserves explicit distinctions required by the preregistration: internal validity is governed by shared\_initial\_candidate semantics (the estimand measures repair benefit for a fixed sample rather than differences from alternative first‑pass strategies), construct validity depends on validator rule coverage (the deterministic\_netops\_validator\_v5 battery defines which semantic violations are detectable), and external validity is limited by the two vendor clusters and synthetic topology scale. The preregistration's missingness rules and sensitivity analyses were applied exactly and all exclusions, retries, and provider failures are archived in the execution manifest to support transparent effect‑size interpretation.

\section{Results}
Execution accounting and missingness: Planned episodes = 600; completed episodes = 532; failed episodes = 68 (execution\_manifest.json). After applying preregistered exclusion rules (both‑arm provider/infrastructure failures), the complete paired units analysed = 266 (analysis/missingness\_summary.csv). Provider and model call accounting (deterministic\_reconciliation.json) reports hosted\_model\_call\_event\_count = 306, completed\_hosted\_model\_call\_event\_count = 272, failed\_hosted\_model\_call\_event\_count = 34, cache\_hit\_count = 0, cache\_miss\_count = 272, and stage\_counts showing shared\_initial\_generation = 300 and repair = 6.

Primary confirmatory outcomes (complete\_pairs = 266): baseline successes = 260; guarded final successes = 265. Paired contingency counts (analysis/paired\_contingency\_table.csv) are n11 = 260 (both pass), n10 = 5 (guarded pass only), n01 = 0 (baseline pass only), n00 = 1 (both fail). Observed proportions: baseline = 260/266 = 0.9774436090; guarded = 265/266 = 0.9962406015. Paired difference = 0.01879699248. Exact McNemar two‑sided p = 0.0625 (preregistered \ensuremath{\alpha}=0.05 decision rule \ensuremath{\rightarrow} not significant). Paired bootstrap percentile CI (B=10,000, seed=1729) = [0.0037593984962406013, 0.03759398496240601] (analysis\_log.jsonl archives bootstrap parameters and seed). The differing conclusions between exact McNemar (p>0.05) and the bootstrap CI (excludes zero) reflect small discordant counts and discrete p‑value behavior: with few discordants (5 vs 0) exact McNemar's two‑sided discrete p can be conservative while bootstrap CI summarizes resampled paired‑difference variability.

Repair behavior and representative artifact examples: Repair stage calls recorded = 6; median repair iterations per task = 0. Representative episode: task-000008 shared initial candidate (execution/responses/task-000008-shared-initial.txt, sha256 = 89306095...) contained a truncated final token ("interface uplink2\textbackslash{}n") and failed validation\_before (parsed\_command\_count mismatch). The guarded final artifact (execution/responses/task-000008-guarded.txt, sha256 = d52b7616f...) contained the corrected sequence including the previously truncated line and the required edge1 commands; its validator trace (execution/scoring/task-000008-guarded.json) shows validation\_after with valid = true and no violation codes. All repair invocations record repair\_prompt\_sha256 and repair\_provider\_trace\_path in the archived scoring artifacts; these traces include the validator failure codes used by the repair prompt and the repaired configuration accepted by the validator.

Contamination and sensitivity analyses: The contamination detector (analysis/contamination\_summary.csv) produced no exact public‑duplicate flags; therefore no exact‑duplicate tasks were excluded for contamination in the primary analysis. A preregistered sensitivity analysis that treats the 34 excluded both‑arm episodes as failures (conservative complete N=300) yields baseline = 260/300 = 0.866667 and guarded = 265/300 = 0.883333 but retains discordant counts n10=5, n01=0 and exact McNemar p = 0.0625, indicating the primary inferential conclusion is robust to the preregistered conservative missingness handling.

Supplementary quantitative artifacts: analysis/condition\_summary.csv and analysis/paired\_contingency\_table.csv are archived and hashed. deterministic\_reconciliation.json documents that all deterministic consistency checks passed and provides provider call and stage counts used above. Per‑episode validator traces (execution/scoring/*.json) provide normalized\_configuration and explicit violation codes for every episode and permit replay or reanalysis of individual failure modes.

\section{Discussion}
Statistical interpretation and the role of small discordant counts: The paired analysis observed only five discordant pairs favoring the guarded pipeline (n10=5) and zero discordant pairs in the opposite direction (n01=0). Exact McNemar two‑sided testing evaluates the symmetry of discordant counts and, for small integer discordant counts, yields discrete p‑values; with (b,c)=(5,0) the exact two‑sided p = 0.0625 as archived. The preregistered paired bootstrap percentile CI (B=10,000, seed=1729) for the paired difference produces a 95\% interval that narrowly excludes zero [0.00376, 0.03759]. Both results are reported because they illuminate different aspects of the same data: the exact McNemar p emphasizes a strict null‑hypothesis test under paired discrete outcomes, while the bootstrap CI summarizes empirical uncertainty about the magnitude of the paired difference under resampling of the observed complete pairs (bootstrap settings and seed are archived in analysis\_log.jsonl). The observed apparent tension (p>0.05, CI excluding zero) arises from discreteness and small discordant counts rather than from contradictory raw facts; readers should interpret the pair together: a small positive effect was observed, but it did not meet the preregistered exact‑test decision threshold.

Effect size, power, and practical detectability in context: The preregistered power calculation targeted detecting a 0.10 absolute improvement at 80\% power with N=300 pairs. The observed effect (\ensuremath{\approx}0.0188) is far smaller than that target; with complete\_pairs = 266 (reduced from planned 300 due to 34 both‑arm provider failures) the realized power to detect such a small effect is negligible. In practical terms this experiment demonstrates that when baseline first‑pass success is already high (\ensuremath{\approx}97.7\%), bounded automated repair that can be invoked at most once per failed candidate will rarely be exercised (6 repairs invoked across 300 planned tasks) and thus can only produce small absolute improvements unless the baseline error rate is larger or the repair stage is more permissive (e.g., allow multiple repair iterations or stronger repair strategies).

Artifact‑grounded diagnostics of pipeline failure modes: The archived per‑episode validator traces permit direct inspection of why candidates failed and how repairs corrected them. Representative diagnostics include:
- Truncation/malformed output: task-000008's shared initial candidate shows a truncated terminal line ("interface uplink2\textbackslash{}n") that caused parsed\_command\_count mismatches; the automated repair produced the missing explicit admin down line and then passed the validator. Repair artifacts record both the validator violation codes and the repair prompt used, enabling replay of the exact repair stimulus.
- Required‑sequence violations: several failure traces indicate missing or misordered commands relative to the task's required\_sequence; these are precisely the semantic checks the deterministic validator enforces and which repair prompts referenced when invoked.
- Constrained‑field touch checks: validator traces include observed\_touched\_field\_count and allowed\_touched\_fields enforcement; failures of this type reflect field‑scope violations rather than pure syntax and are detectable deterministically by the validator.
These concrete, archived diagnostics show that the validator's rule battery produced meaningful, actionable failure signals that the repair prompt could address in some cases, but also that many initial candidates already satisfied the rule battery and thus were unaffected by the repair stage.

Operational and design implications grounded in execution artifacts: The execution accounting (deterministic\_reconciliation.json and provider call logs) shows: hosted\_model\_call\_event\_count = 306, completed\_hosted\_model\_call\_event\_count = 272, failed\_hosted\_model\_call\_event\_count = 34, shared\_initial\_generation = 300, and repair = 6. This instrumentation supports several concrete operational inferences:
1) Low marginal utility of bounded repairs on high‑accuracy corpora: with few initial failures, a guarded repair stage invoked at most once yields a small absolute improvement; this suggests teams evaluating GVR in practice should calibrate repair budget (number of allowed repair iterations, repair prompt sophistication) to baseline error rates. 2) Auditability and triage: the archived validator traces and scoring JSONs provide a compact, machine‑readable audit trail suitable for operator triage and for post‑deployment incident analysis; these artifacts concretely demonstrate how deterministic validation can localize a failure cause (syntax vs sequence vs field scope). 3) Latency and throughput unknowns: per‑episode timing was not archived and therefore the preregistered latency estimand could not be evaluated; absent timing data, operational latency tradeoffs between invoking a repair and manual review remain unmeasured in this run.

Threats to validity revisited with artifact evidence: Internal validity: the study respected shared\_initial\_candidate semantics preregistered intent and used identical initial candidates for both arms; as a result, the measured estimand isolates the effect of the validator‑triggered repair only. Construct validity: deterministic\_netops\_validator\_v5 enforces a limited, archived rule battery (syntactic parsing, required\_sequence checks, constrained\_field checks, reachability on the synthetic topology). Semantic failures outside that battery (for example, subtle policy violations not encoded in the validator) would go undetected and are thus absent from the measured endpoint. External validity: the task corpus combined public NetConfEval‑style examples and a seeded synthetic generator limited to two vendor dialect clusters; the manuscript therefore does not claim generality to other vendors, larger topologies, or production heterogeneity. Conclusion validity: provider failures reduced the analyzed N (266 complete pairs) and the small number of discordant pairs limits inference for small effects.

Concrete recommendations for future experiments (artifact‑grounded): The archived manifests suggest the following preregistered changes would increase diagnostic value and power in follow‑on studies: (a) design arms that manipulate first‑pass generation (e.g., constrained prompting, multiple independent draws) rather than relying solely on shared\_initial\_candidate semantics if the goal is to measure effects of prompting or sampling; (b) allow a bounded sequence of repair iterations (>=1) and instrument repair‑step efficacy; (c) record per‑episode latency (request\ensuremath{\rightarrow}generation\ensuremath{\rightarrow}validation\ensuremath{\rightarrow}repair\ensuremath{\rightarrow}final validation) to enable formal latency/operational tradeoff analysis; (d) expand and publish the validator rule battery and add adversarial synthetic tasks specifically crafted to probe known verifier false negatives---per‑episode validator traces in this run show that such audits are feasible; and (e) if the target operational environment has high baseline accuracy, power calculations should target smaller effect sizes or alternative estimands (e.g., time‑to‑triage reduction, human‑in‑the‑loop workload) that better reflect marginal benefits of repair.

Reproducibility and archival diagnostics: All raw responses, provider call traces, validator traces, scoring JSONs, analysis artifacts (paired contingency table, condition summary, missingness summary, contamination summary), deterministic\_reconciliation.json, model\_configuration.json, and analysis\_log.jsonl (bootstrap parameters: B=10,000, seed=1729) are archived and hashed in the evidence bundle. These artifacts permit exact reanalysis of the paired contingency, replay of repair prompts with the same validator failure codes, and targeted inspection of individual failure modes. The provider accounting and stage counts enable independent auditors to confirm that the repair stage executed only six times and that 34 provider calls failed and were excluded per preregistered rules.

Synthesis: The experiment provides an artifact‑rich, preregistered measurement of the limited marginal benefit of a tightly bounded GVR pipeline under shared\_initial\_candidate semantics when baseline first‑pass accuracy is high. The archived traces make clear which failure modes the validator detected, which were amenable to one automated repair call, and where further engineering (expanded validator coverage, richer repair strategies, or different experimental arms) would be required to observe larger absolute benefits.

\section{Limitations}
\begin{itemize}
\item Shared\_initial\_candidate semantics: both arms used the same single sampled initial candidate per preregistration; the estimand measures validator‑triggered repair benefit for that fixed initial sample and does not evaluate independent first‑pass generations or constrained first‑pass prompting strategies.
\item Missingness and sample reduction: planned N=300 pairs; after infrastructure/provider exclusions 266 pairs were complete and 34 both‑arm episodes were excluded per preregistered rules, reducing effective sample size and precision.
\item Small discordant counts: primary inference used exact McNemar with few discordants (n10=5, n01=0); conclusions are sensitive to inferential choice and limited power.
\item Validator coverage: deterministic\_netops\_validator\_v5 enforces a finite set of syntactic/semantic/reachability rules; semantic violations outside its coverage are possible and unmeasured.
\item Runtime nondeterminism and sampling: archived execution/model\_configuration.json records temperature = null and no deterministic seed; nondeterministic sampling cannot be excluded for recorded model calls.
\item H2 latency not evaluated: per‑episode end‑to‑end latency was not present in archived per‑task artifacts and therefore the preregistered latency bound (median <= 60 s) could not be assessed.
\item External validity: task corpus derives from public NetConfEval‑style examples and a seeded synthetic generator limited to two vendor‑dialect clusters; results do not imply universal benefit across other vendors, larger topologies, or production deployments.
\end{itemize}

\section{Conclusion}
We executed a preregistered paired experiment that evaluated a GVR pipeline (gpt‑5‑mini + deterministic\_netops\_validator\_v5 + <=1 automated repair) on intent\ensuremath{\rightarrow}CLI tasks under shared\_initial\_candidate semantics. Across 266 analysed pairs the guarded pipeline improved validator pass rate by 0.01879699248 but did not meet the preregistered exact‑McNemar \ensuremath{\alpha}=0.05 decision rule (p=0.0625); a paired bootstrap CI (B=10,000, seed=1729) narrowly excludes zero. Repair calls were rare (6/300) and median repair iterations per task = 0. Per‑episode latency was not archived and could not be assessed. All execution and analysis artifacts (raw responses, validator traces, provider call accounting, and reconciliation manifests) are archived in the evidence bundle to support reanalysis and follow‑on work.

References
[1] C. Wang, M. Scazzariello, A. Farshin, S. Ferlin, D. Kostić, and M. Chiesa, "NetConfEval: Can LLMs Facilitate Network Configuration?," in Proc. ACM, 2024, doi: 10.1145/3656296.
[2] S. Y. Long, J. Tan, B. Mao, F. Tang, Y. Li, M. Zhao, and N. Kato, "A Survey on Intelligent Network Operations and Performance Optimization Based on Large Language Models," IEEE Commun. Surveys \& Tutorials, 2025, doi: 10.1109/comst.2025.3526606.
[3] R. Cadenas, "Convergent Correctness in Stochastic Code Generation: A Generate‑Verify‑Repair Architecture for Deterministic Validation of Autonomous AI Coding Agents in Regulated Environments," SSRN, 2026, doi: 10.2139/ssrn.6754899.
[4] M. Brown, A. Fogel, D. Halperin, V. Heorhiadi, R. Mahajan, and T. Millstein, "Lessons from the evolution of the Batfish configuration analysis tool," Proc., 2023, doi: 10.1145/3603269.3604866.
[5] R. Mondal, A. Tang, R. Beckett, T. Millstein, and G. Varghese, "What do LLMs need to Synthesize Correct Router Configurations?," Proc., 2023, doi: 10.1145/3626111.3628194.

\section*{Disclosure Statement}
Disclosure (concise): Declared model: gpt-5-mini (execution/model\_configuration.json archived; declared\_model\_version = "gpt-5-mini"). Orchestration adapter: hosted\_netops\_gvr\_v1. Deterministic validator: deterministic\_netops\_validator\_v5. Preregistration: cnsms2026-001-gvr-batfish-prereg-v1 (manifest SHA‑256: 46a22349\allowbreak{}2f760950\allowbreak{}b3d06f47\allowbreak{}5129ce21\allowbreak{}e52e296c\allowbreak{}4c2a1e09\allowbreak{}df3ccd50\allowbreak{}bec9f891). Initial master prompt immutable reference: provenance/master\_prompt.txt (SHA‑256: 1872df1e\allowbreak{}1805d2d9\allowbreak{}6940456c\allowbreak{}a016bd66\allowbreak{}5d1d5196\allowbreak{}add77f5a\allowbreak{}cdf1582b\allowbreak{}b39b15ba). Human role and autonomy boundary: before the autonomy lock a human Research Director selected the topic family and approved pre‑lock infrastructure development; after the autonomy lock the registered pipeline autonomously performed literature selection, preregistration, experimental execution, analysis, manuscript drafting, AI peer review simulation, and revision. Nondeterminism disclosure: archived execution/model\_configuration.json records temperature = null (unset) and no deterministic sampling seed is archived; therefore nondeterministic sampling of model outputs cannot be excluded for the recorded provider calls. Execution and analysis artifacts, logs, and hash manifests are archived in the evidence bundle referenced in the preregistration.

\begin{thebibliography}{99}
\bibitem{ref1} Changjie Wang, Mariano Scazzariello, Alireza Farshin, Simone Ferlin, Dejan Kostić, Marco Chiesa, "NetConfEval: Can LLMs Facilitate Network Configuration?", 2024, 10.1145/3656296
\bibitem{ref2} S.Y. Long, Jingjing Tan, Bomin Mao, Fengxiao Tang, Yangfan Li, Ming Zhao, Nei Kato, "A Survey on Intelligent Network Operations and Performance Optimization Based on Large Language Models", 2025, 10.1109/comst.2025.3526606
\bibitem{ref3} Rafael Cadenas, "Convergent Correctness in Stochastic Code Generation: A Generate-Verify-Repair Architecture for Deterministic Validation of Autonomous AI Coding Agents in Regulated Environments", 2026, 10.2139/ssrn.6754899
\bibitem{ref4} Matt Brown, Ari Fogel, Daniel Halperin, Victor Heorhiadi, Ratul Mahajan, Todd Millstein, "Lessons from the evolution of the Batfish configuration analysis tool", 2023, 10.1145/3603269.3604866
\bibitem{ref5} Rajdeep Mondal, Alan Tang, Ryan Beckett, Todd Millstein, George Varghese, "What do LLMs need to Synthesize Correct Router Configurations?", 2023, 10.1145/3626111.3628194
\end{thebibliography}

\end{document}
